% TOP_PROBLEM: {"id": 30005155, "title": "Fractal Schrödinger Convergence Above Half Dimension", "category_id": 8, "category": {"id": 8, "name": "analysis", "display_name": "Analysis", "description": "Limits, continuity, calculus, and function theory.", "slug": "analysis", "order_index": 8, "created_at": "2026-07-31T15:26:25.671Z"}, "status": "open"}
% FIRST_POSTED: 2026-10-06
% ATTEMPT_STATUS: unresolved
% OPENAI_CLAIM: {"schema_version":1,"source_repo":"https://github.com/openai/math","source_commit":"adc7f1241b42e322a6451854ab7e4b4c146bf78a","release_date":"2026-10-06","match":"partial","resolution":"partial","independently_reviewed":false,"second_pass":"n/a","families":[{"family":"080","title":"The exact Sobolev endpoint for Schrödinger convergence","subject":"Real and complex analysis","manuscripts":[{"title":"Endpoint convergence for the planar Schrödinger equation","pdf_path":"preprints/Endpoint-convergence-for-the-planar-Schrodinger-equation-September-24-2026/paper.pdf","readme_path":"preprints/Endpoint-convergence-for-the-planar-Schrodinger-equation-September-24-2026/README.md","date":"2026-09-24","theorem_ref":"thm:main (sections/01-introduction.tex:23)","statement_sources":[{"path":"preprints/Endpoint-convergence-for-the-planar-Schrodinger-equation-September-24-2026/build/sections/01-introduction.tex","line":23,"label":"thm:main"}]},{"title":"Endpoint pointwise convergence for the Schrödinger equation in higher dimensions","pdf_path":"preprints/Endpoint-pointwise-convergence-for-the-Schrodinger-equation-in-higher-dimensions-September-24-2026/paper.pdf","readme_path":"preprints/Endpoint-pointwise-convergence-for-the-Schrodinger-equation-in-higher-dimensions-September-24-2026/README.md","date":"2026-09-24","theorem_ref":"m:main (sections/introduction.tex:34)","statement_sources":[{"path":"preprints/Endpoint-pointwise-convergence-for-the-Schrodinger-equation-in-higher-dimensions-September-24-2026/build/sections/introduction.tex","line":34,"label":"m:main"}]}],"lean":null,"reasoning_trace":null}]}
\documentclass[11pt]{article}
\usepackage[margin=25mm]{geometry}
\usepackage{fontspec}
\setmainfont{Latin Modern Roman}
\usepackage{amsmath,amssymb,mathtools}
\usepackage{unicode-math}
\setmathfont{Latin Modern Math}
\usepackage{enumitem}
\usepackage{xurl,hyperref,bookmark}
\hypersetup{colorlinks=true,urlcolor=blue}
\setcounter{secnumdepth}{0}
\setlength{\parindent}{0pt}
\setlength{\parskip}{5pt}
\setlength{\emergencystretch}{3em}
\begin{document}
\section*{Fractal Schrödinger Convergence Above Half Dimension}

\subsection{OpenAI's claim}
\textit{Planar Schrödinger endpoint}: The paper claims almost-everywhere convergence to the initial datum for planar free Schrödinger evolution at Sobolev order $1/3$; gaussian regularization is removed on a common full-measure spatial set before taking the time limit.

\textit{Higher-dimensional Schrödinger endpoints}: For every $n\ge3$, the claimed critical Sobolev order for almost-everywhere free Schrödinger convergence is $n/(2(n+1))$; the regularized evolution is asserted to have a common spatial representative throughout the positive time interval.

\subsection{Mathematical statement of the claimed result}
\paragraph{Mathematical claim 1.} In spatial dimension $2$ put $s=1/3$ and $U_af(x,t)=(2\pi)^{-n}\int_{\mathbb R^n}e^{ix\cdot\xi-it|\xi|^2-a|\xi|^2}\widehat f(\xi)\,d\xi$, using $n=2$ in the planar case. For every complex $f\in H^s(\mathbb R^n)$ the claim provides one full-measure set $E_f$ of Lebesgue points and a function $v$ continuous in $t\in[0,1)$, with $v(x,0)=f^*(x)$, such that \[\lim_{a\downarrow0}\sup_{0<t<1}|U_af(x,t)-v(x,t)|=0\quad(x\in E_f).\] Consequently the Gaussian-regularized Schr\"odinger evolution tends to the Lebesgue-point representative $f^*$ as $t\downarrow0$. The exceptional spatial set is independent of time.

\paragraph{Mathematical claim 2.} In spatial dimension $n\ge3$ put $s=n/(2(n+1))$ and $U_af(x,t)=(2\pi)^{-n}\int_{\mathbb R^n}e^{ix\cdot\xi-it|\xi|^2-a|\xi|^2}\widehat f(\xi)\,d\xi$. For every complex $f\in H^s(\mathbb R^n)$ the claim provides one full-measure set $E_f$ of Lebesgue points and a function $v$ continuous in $t\in[0,1)$, with $v(x,0)=f^*(x)$, such that \[\lim_{a\downarrow0}\sup_{0<t<1}|U_af(x,t)-v(x,t)|=0\quad(x\in E_f).\] Consequently the Gaussian-regularized Schr\"odinger evolution tends to the Lebesgue-point representative $f^*$ as $t\downarrow0$. The exceptional spatial set is independent of time.

\subsection{Outline of the argument}
\textit{Planar Schrödinger endpoint}: At unit frequency, a spatial measure on a measurable time graph is analyzed using transverse packets, sparse configurations, and a separate frame estimate. Flat-frequency and geometric reductions give a scale-uniform dyadic distribution bound. Endpoint summation and uniform approximation construct one representative of the evolution with the asserted limiting behavior.

\textit{Higher-dimensional Schrödinger endpoints}: Transverse packet fields are controlled by a frame count with polynomial dependence on transversality, combined with amplitude and occupancy hierarchies. Induction produces a short-time maximal estimate at precisely the critical frequency factor. A single time tree is used across frequencies for square summation, after which Gaussian approximation removes exceptional sets uniformly in time.

\subsection{Scope relative to this problem}
\textit{Planar Schrödinger endpoint}: The endpoint theorem treats Lebesgue-almost-everywhere convergence, the alpha=d case of the fractal convergence question, in its stated planar or n>=3 dimensions. It does not determine the sharp s(alpha) for d/2<alpha<d.

\textit{Higher-dimensional Schrödinger endpoints}: The endpoint theorem treats Lebesgue-almost-everywhere convergence, the alpha=d case of the fractal convergence question, in its stated planar or n>=3 dimensions. It does not determine the sharp s(alpha) for d/2<alpha<d.

This page records partial support. Complete resolution of the catalogue question is not established by this record.

\subsection{Formal verification}
\textit{Planar Schrödinger endpoint}: The pinned catalogue lists no Lean formalization for this result.

\textit{Higher-dimensional Schrödinger endpoints}: The pinned catalogue lists no Lean formalization for this result.

This site has not run the Lean development. These links describe OpenAI’s selected formal statements and should be read with their stated scope.

\subsection{Status caveat}
OpenAI reports that an internal model produced the manuscripts. Their arguments have not been independently verified as proofs by this site. Any portion outside the documented formal scope may contain mistakes, and the catalogue label reports the announced mathematical claim.

\subsection{Sources}
\begin{itemize}
\item OpenAI, \emph{Endpoint convergence for the planar Schrödinger equation}, 2026-09-24; citation key \texttt{OAI:Endpoint-convergence-for-the-planar-Schrodinger-equation-September-24-2026}. \href{https://github.com/openai/math/blob/adc7f1241b42e322a6451854ab7e4b4c146bf78a/preprints/Endpoint-convergence-for-the-planar-Schrodinger-equation-September-24-2026/paper.pdf}{PDF (GitHub)}; \href{https://github.com/openai/math/raw/adc7f1241b42e322a6451854ab7e4b4c146bf78a/preprints/Endpoint-convergence-for-the-planar-Schrodinger-equation-September-24-2026/paper.pdf}{PDF (direct)}; \href{https://github.com/openai/math/blob/adc7f1241b42e322a6451854ab7e4b4c146bf78a/preprints/Endpoint-convergence-for-the-planar-Schrodinger-equation-September-24-2026/README.md}{Manuscript page and citation}; \href{https://github.com/openai/math/blob/main/preprints/Endpoint-convergence-for-the-planar-Schrodinger-equation-September-24-2026/paper.pdf}{Latest version}; Statement source: \href{https://github.com/openai/math/blob/adc7f1241b42e322a6451854ab7e4b4c146bf78a/preprints/Endpoint-convergence-for-the-planar-Schrodinger-equation-September-24-2026/build/sections/01-introduction.tex\#L23}{thm:main}.
\item OpenAI, \emph{Endpoint pointwise convergence for the Schrödinger equation in higher dimensions}, 2026-09-24; citation key \texttt{OAI:Endpoint-pointwise-convergence-for-the-Schrodinger-equation-in-higher-dimensions-September-24-2026}. \href{https://github.com/openai/math/blob/adc7f1241b42e322a6451854ab7e4b4c146bf78a/preprints/Endpoint-pointwise-convergence-for-the-Schrodinger-equation-in-higher-dimensions-September-24-2026/paper.pdf}{PDF (GitHub)}; \href{https://github.com/openai/math/raw/adc7f1241b42e322a6451854ab7e4b4c146bf78a/preprints/Endpoint-pointwise-convergence-for-the-Schrodinger-equation-in-higher-dimensions-September-24-2026/paper.pdf}{PDF (direct)}; \href{https://github.com/openai/math/blob/adc7f1241b42e322a6451854ab7e4b4c146bf78a/preprints/Endpoint-pointwise-convergence-for-the-Schrodinger-equation-in-higher-dimensions-September-24-2026/README.md}{Manuscript page and citation}; \href{https://github.com/openai/math/blob/main/preprints/Endpoint-pointwise-convergence-for-the-Schrodinger-equation-in-higher-dimensions-September-24-2026/paper.pdf}{Latest version}; Statement source: \href{https://github.com/openai/math/blob/adc7f1241b42e322a6451854ab7e4b4c146bf78a/preprints/Endpoint-pointwise-convergence-for-the-Schrodinger-equation-in-higher-dimensions-September-24-2026/build/sections/introduction.tex\#L34}{m:main}.
\item \href{https://github.com/openai/math/blob/adc7f1241b42e322a6451854ab7e4b4c146bf78a/CONTENTS.md}{OpenAI catalogue entry}.
\end{itemize}
\end{document}
